% written by gen/mk.py: (a) the eight symmetries of the square by coset of {identity, hflip}; numbers: data/mv2.json (lowest of the five ID data sets, train images).  (b), (c): schematics.
\begin{tikzpicture}[x=1mm, y=1mm, font=\figlab, line cap=round, line join=round, >={Stealth[length=1.6mm, width=1.3mm]}]
\node[anchor=south west, font=\scriptsize, inner sep=0pt] at (-1,22.5) {\textbf{(a)}\ 正方形の8つの対称変換の中の転置};
\begin{scope}[shift={(5.0,9)}]\draw[black!45, line width=0.4pt] (-4.3,-4.3) rectangle (4.3,4.3);\begin{scope}[cm={1,0,0,1,(0,0)}]\draw[black, line width=0.9pt] (-3.4,-3.1) -- (3.4,-3.1);\draw[black, line width=0.7pt] (-1.2,-3.1) -- (-1.2,3.1);\fill[redc] (-1.2,3.1) -- (2.5,2.0) -- (-1.2,0.9) -- cycle;\end{scope}\end{scope}
\node[anchor=north, inner sep=1pt] at (5.0,4.1) {そのまま};
\node[anchor=north, inner sep=1pt] at (5.0,0.6) {--};
\begin{scope}[shift={(15.2,9)}]\draw[black!45, line width=0.4pt] (-4.3,-4.3) rectangle (4.3,4.3);\begin{scope}[cm={-1,0,0,1,(0,0)}]\draw[black, line width=0.9pt] (-3.4,-3.1) -- (3.4,-3.1);\draw[black, line width=0.7pt] (-1.2,-3.1) -- (-1.2,3.1);\fill[redc] (-1.2,3.1) -- (2.5,2.0) -- (-1.2,0.9) -- cycle;\end{scope}\end{scope}
\node[anchor=north, inner sep=1pt] at (15.2,4.1) {左右反転};
\node[anchor=north, inner sep=1pt] at (15.2,0.6) {49.5};
\draw[black!70, line width=0.5pt] (0.7000000000000002,15.6) -- (0.7000000000000002,16.6) -- (19.5,16.6) -- (19.5,15.6);
\node[anchor=south, text=black!70, inner sep=1.2pt] at (10.1,16.6) {分布内にとどまる};
\begin{scope}[shift={(28.0,9)}]\draw[black!45, line width=0.4pt] (-4.3,-4.3) rectangle (4.3,4.3);\begin{scope}[cm={-1,0,0,-1,(0,0)}]\draw[black, line width=0.9pt] (-3.4,-3.1) -- (3.4,-3.1);\draw[black, line width=0.7pt] (-1.2,-3.1) -- (-1.2,3.1);\fill[redc] (-1.2,3.1) -- (2.5,2.0) -- (-1.2,0.9) -- cycle;\end{scope}\end{scope}
\node[anchor=north, inner sep=1pt] at (28.0,4.1) {180度回転};
\node[anchor=north, inner sep=1pt] at (28.0,0.6) {63.5};
\begin{scope}[shift={(38.2,9)}]\draw[black!45, line width=0.4pt] (-4.3,-4.3) rectangle (4.3,4.3);\begin{scope}[cm={1,0,0,-1,(0,0)}]\draw[black, line width=0.9pt] (-3.4,-3.1) -- (3.4,-3.1);\draw[black, line width=0.7pt] (-1.2,-3.1) -- (-1.2,3.1);\fill[redc] (-1.2,3.1) -- (2.5,2.0) -- (-1.2,0.9) -- cycle;\end{scope}\end{scope}
\node[anchor=north, inner sep=1pt] at (38.2,4.1) {上下反転};
\node[anchor=north, inner sep=1pt] at (38.2,0.6) {61.9};
\draw[black!70, line width=0.5pt] (23.7,15.6) -- (23.7,16.6) -- (42.5,16.6) -- (42.5,15.6);
\node[anchor=south, text=black!70, inner sep=1.2pt] at (33.1,16.6) {上下が逆};
\begin{scope}[shift={(51.00000000000001,9)}]\draw[black!45, line width=0.4pt] (-4.3,-4.3) rectangle (4.3,4.3);\begin{scope}[cm={0,1,-1,0,(0,0)}]\draw[black, line width=0.9pt] (-3.4,-3.1) -- (3.4,-3.1);\draw[black, line width=0.7pt] (-1.2,-3.1) -- (-1.2,3.1);\fill[redc] (-1.2,3.1) -- (2.5,2.0) -- (-1.2,0.9) -- cycle;\end{scope}\end{scope}
\node[anchor=north, inner sep=1pt] at (51.00000000000001,4.1) {90度回転};
\node[anchor=north, inner sep=1pt] at (51.00000000000001,0.6) {76.3};
\begin{scope}[shift={(61.2,9)}]\draw[black!45, line width=0.4pt] (-4.3,-4.3) rectangle (4.3,4.3);\begin{scope}[cm={0,1,1,0,(0,0)}]\draw[black, line width=0.9pt] (-3.4,-3.1) -- (3.4,-3.1);\draw[black, line width=0.7pt] (-1.2,-3.1) -- (-1.2,3.1);\fill[redc] (-1.2,3.1) -- (2.5,2.0) -- (-1.2,0.9) -- cycle;\end{scope}\end{scope}
\node[anchor=north, inner sep=1pt] at (61.2,4.1) {反転置};
\node[anchor=north, inner sep=1pt] at (61.2,0.6) {76.9};
\draw[navyc, line width=0.5pt] (46.70000000000001,15.6) -- (46.70000000000001,16.6) -- (65.5,16.6) -- (65.5,15.6);
\node[anchor=south, text=navyc, inner sep=1.2pt] at (56.1,16.6) {横倒し};
\begin{scope}[shift={(74.0,9)}]\draw[black!45, line width=0.4pt] (-4.3,-4.3) rectangle (4.3,4.3);\begin{scope}[cm={0,-1,1,0,(0,0)}]\draw[black, line width=0.9pt] (-3.4,-3.1) -- (3.4,-3.1);\draw[black, line width=0.7pt] (-1.2,-3.1) -- (-1.2,3.1);\fill[redc] (-1.2,3.1) -- (2.5,2.0) -- (-1.2,0.9) -- cycle;\end{scope}\end{scope}
\node[anchor=north, inner sep=1pt] at (74.0,4.1) {270度回転};
\node[anchor=north, inner sep=1pt] at (74.0,0.6) {76.8};
\begin{scope}[shift={(84.2,9)}]\draw[redc, line width=0.9pt] (-4.3,-4.3) rectangle (4.3,4.3);\begin{scope}[cm={0,-1,-1,0,(0,0)}]\draw[black, line width=0.9pt] (-3.4,-3.1) -- (3.4,-3.1);\draw[black, line width=0.7pt] (-1.2,-3.1) -- (-1.2,3.1);\fill[redc] (-1.2,3.1) -- (2.5,2.0) -- (-1.2,0.9) -- cycle;\end{scope}\end{scope}
\node[anchor=north, inner sep=1pt, text=redc] at (84.2,4.1) {転置};
\node[anchor=north, inner sep=1pt] at (84.2,0.6) {76.6};
\draw[navyc, line width=0.5pt] (69.7,15.6) -- (69.7,16.6) -- (88.5,16.6) -- (88.5,15.6);
\node[anchor=south, text=navyc, inner sep=1.2pt] at (79.10000000000001,16.6) {横倒し};
\node[anchor=north west, inner sep=1pt, align=left, text width=88mm] at (-1,-3.2) {下の数：その変換をした train 画像を正立の train 画像と見分ける AUROC(\%)の, 5つの分布内での最小. 横倒しにする4つだけが, どの分布内でも見分けられる.};
\begin{scope}[shift={(99.5,0)}]
\node[anchor=south west, font=\scriptsize, inner sep=0pt] at (-1,22.5) {\textbf{(b)}\ 記述子の2つの成分};
\coordinate (G) at (4,3); \coordinate (GT) at (30,15); \coordinate (E) at (17,9);
\draw[black!55, line width=0.5pt] (G) -- (GT); \draw[->, navyc, line width=0.9pt] (E) -- (G); \fill[navyc] (G) circle (0.75); \draw[navyc, fill=white, line width=0.6pt] ($(GT)+(-0.7,-0.7)$) rectangle ++(1.4,1.4); \fill[black] (E) circle (0.5);
\node[anchor=north, text=navyc, inner sep=2pt] at ($(G)+(0,-0.6)$) {$g(x)$}; \node[anchor=south, text=navyc, inner sep=2pt] at ($(GT)+(0,0.8)$) {$g(\Tm x)$}; \node[anchor=south east, inner sep=1.5pt] at ($(E)+(0.2,0.5)$) {$e(x)$};
\node[anchor=north west, text=navyc, inner sep=1pt] at ($(G)!0.5!(E)+(0.4,-0.6)$) {$o(x)$};
\draw[black, line width=0.4pt, |<->|] ($(G)+(-2.2,4.77)$) -- node[above, sloped, inner sep=1.5pt] {$\tau=2|o(x)|$} ($(GT)+(-2.2,4.77)$);
\node[anchor=north west, inner sep=1pt, align=left, text width=36mm] at (-1,-3.2) {$e$：転置で変わらない成分.\\$o$：転置で符号が変わる成分.};
\end{scope}
\begin{scope}[shift={(140.5,0)}]
\node[anchor=south west, font=\scriptsize, inner sep=0pt] at (-1,22.5) {\textbf{(c)}\ 定理1：ずれは対数密度の勾配の差};
\foreach \a/\b in {8/5, 5.2/3.2, 2.5/1.5}{\draw[navyc, line width=0.5pt, rotate around={20:(7.5,1.5)}] (7.5,1.5) ellipse (\a mm and \b mm);}
\foreach \a/\b in {8/5, 5.2/3.2, 2.5/1.5}{\draw[redc, densely dashed, line width=0.5pt, rotate around={70:(33,1.5)}] (33,1.5) ellipse (\a mm and \b mm);}
\node[text=navyc, anchor=north, inner sep=1pt] at (7.5,-4.6) {正立の分布}; \node[text=redc, anchor=north, inner sep=1pt] at (33,-6.9) {転置した分布};
\coordinate (Y) at (20,20.4); \coordinate (A) at (14.6,13.2); \coordinate (Bq) at (25.4,13.2);
\draw[->, navyc, line width=0.8pt] (Y) -- (A); \draw[->, redc, line width=0.8pt] (Y) -- (Bq); \draw[->, black, line width=1.0pt] (A) -- (Bq); \fill[black] (Y) circle (0.6);
\node[anchor=west, inner sep=2pt] at ($(Y)+(0.6,0.4)$) {入力 $y$}; \node[anchor=east, text=navyc, inner sep=1.5pt] at ($(A)+(-0.6,0)$) {$D(y)$}; \node[anchor=west, text=redc, inner sep=1.5pt] at ($(Bq)+(0.6,0)$) {$\Tm D(\Tm y)$};
\node[anchor=north, inner sep=2pt] at (20,12.6) {ずれ $\Delta(y)$};
\end{scope}
\end{tikzpicture}
